\documentclass[11pt]{article}

\usepackage[margin=1in]{geometry}
\usepackage{amsmath,amssymb}
\usepackage[T1]{fontenc}
\usepackage[utf8]{inputenc}
\usepackage{lmodern}
\usepackage{graphicx}
\usepackage{booktabs}
\usepackage{float}
\usepackage{hyperref}

\newcommand{\uc}{\ensuremath{\mu\text{C}/\text{cm}^2}}
\newcommand{\mvcm}{\ensuremath{\text{MV}/\text{cm}}}

\title{Physics and Mathematical Formulation of the Active Multidomain Preisach Hysteresis Model}
\author{Krishna Chemudupati}
\date{\today}

\begin{document}
\maketitle

\begin{abstract}
This note records the exact physics and mathematical structure used by the
active multidomain Preisach hysteresis model in the \texttt{src/} tree. The
purpose is to state precisely what is being solved in the current repository,
rather than to survey the full range of possible ferroelectric models. The
active workflow combines a multidomain Preisach constitutive law for AlScN,
one-dimensional stack electrostatics with finite metal screening and a
work-function-derived built-in term, and a bisection-based quasistatic
self-consistency solver. The equations, parameters, and numerical protocol are
written here in the same form used by the active implementation.
\end{abstract}

\tableofcontents
\clearpage

\section{Problem statement}

The active repository solves a quasistatic polarization--voltage hysteresis
problem for stacks of the form
\[
\text{top metal}
\;|\;
\text{interlayer}
\;|\;
\text{AlScN}
\;|\;
\text{dead layer}
\;|\;
\text{bottom metal}.
\]
In the main active figures the dead layer is set to zero thickness, but the
stack definition retains the dead-layer terms symbolically because they remain
part of the implemented electrostatic model. The question answered by the
active workflow is therefore narrow and concrete: for a fixed multilayer stack,
a fixed domain population, and a prescribed voltage sweep, what self-consistent
multidomain ferroelectric polarization is obtained at each bias point?

\section{Variables and units}

The macroscopic stack variables are the interlayer thickness $t_i$, the
ferroelectric thickness $t_f$, the dead-layer thickness $t_d$, the dielectric
constants $k_i$, $k_f$, and $k_d$, the top and bottom screening lengths
$\lambda_t$ and $\lambda_b$, the corresponding screening dielectric factors
$k_t$ and $k_b$, the top and bottom work functions $W_t$ and $W_b$, the
applied voltage $V_a$, the screening charge density $\sigma_s$, the dead-layer
polarization $P_d$, and the net ferroelectric polarization
$P_{\mathrm{FE}}$.

The ferroelectric is represented as a population of $N$ Preisach domains.
Domain $n$ is described by a structural parameter $\eta_n = (c/a)_n$, a local
saturation polarization $P_{s,n}$, a local coercive field $E_{c,n}$, and a
binary state variable $s_n \in \{-1,+1\}$. The net ferroelectric polarization
is defined as
\begin{equation}
P_{\mathrm{FE}} = \frac{1}{N}\sum_{n=1}^{N} s_n P_{s,n}.
\label{eq:pfe-average-mifm}
\end{equation}

Internally, the code works in atomic units. This is handled in
\texttt{src/atomicunits.py}, where voltage is converted from volts to atomic
units, thickness is converted from nanometers to bohr, polarization is
converted from \uc\ to atomic units, and electric field is converted from
\mvcm\ to atomic units. The electrostatic constant used internally is
\begin{equation}
\epsilon_0^{(\mathrm{au})} = \frac{1}{4\pi}.
\label{eq:eps0-au-mifm}
\end{equation}
All equations below are written in a form consistent with that implementation.

\section{Multidomain Preisach constitutive law}

The active model generates the ferroelectric domain population from a
structural-disorder proxy, namely the local lattice ratio $c/a$. For each
domain,
\begin{equation}
\eta_n \sim \mathcal{N}(\mu_{c/a}, \sigma_{c/a}).
\label{eq:ca-dist-mifm}
\end{equation}
In the active figure workflow the default values are
\begin{equation}
\mu_{c/a} = 1.54,
\qquad
\sigma_{c/a} = 0.04.
\label{eq:ca-active-mifm}
\end{equation}

If no override values are supplied, the local saturation polarization and the
local coercive field are generated through the default constitutive maps
\begin{align}
P_{s,n} &= 333.33\,\eta_n - 400
\qquad \text{(\uc, before unit conversion)}, \label{eq:ps-map-mifm} \\
E_{c,n} &= 3.16\,\eta_n - 1.1
\qquad \text{(\mvcm, before unit conversion)}. \label{eq:ec-map-mifm}
\end{align}
The resulting values are then converted into atomic units. The class also
supports an override path in which user-specified means and standard deviations
for $P_s$ and $E_c$ are imposed by constructing linear maps from the sampled
$\eta_n$ values,
\begin{align}
P_{s,n} &= \left| a_P \eta_n + b_P \right|, \\
E_{c,n} &= \left| a_E \eta_n + b_E \right|,
\end{align}
with
\begin{align}
a_P &= \frac{\sigma_{P_s}}{\sigma_{c/a}},
&
b_P &= \mu_{P_s} - a_P \mu_{c/a}, \\
a_E &= \frac{\sigma_{E_c}}{\sigma_{c/a}},
&
b_E &= \mu_{E_c} - a_E \mu_{c/a}.
\end{align}
The absolute value in this override path is used only to prevent unphysical
negative local magnitudes.

The switching law itself is a symmetric rectangular hysteron. If $E$ denotes a
candidate internal field seen by the ferroelectric, the domain update is
\begin{equation}
s_n(E) =
\begin{cases}
-1, & s_n^{\mathrm{base}} = +1 \text{ and } E \le -E_{c,n}, \\
+1, & s_n^{\mathrm{base}} = -1 \text{ and } E \ge +E_{c,n}, \\
s_n^{\mathrm{base}}, & \text{otherwise}.
\end{cases}
\label{eq:switch-rule-mifm}
\end{equation}
The current active model therefore assumes binary local states, abrupt
threshold-driven switching, and no continuous local $P(E)$ law or explicit
switching kinetics.

\section{Stack electrostatics}

The stack is instantiated by the \texttt{FerroelectricDiode} class and carries
the geometry and boundary-condition data required by the electrostatic solver.
The active implementation uses the thicknesses $t_i$, $t_f$, and $t_d$, the
dielectric constants $k_i$, $k_f$, and $k_d$, the screening constants
$k_t$ and $k_b$, the screening lengths $\lambda_t$ and $\lambda_b$, the work
functions $W_t$ and $W_b$, and the current polarization
$P_{\mathrm{FE}}$. In the current code the dead-layer parameters are
\begin{equation}
k_d = \frac{k_f}{2},
\qquad
P_d = 0.
\label{eq:dead-layer-active-mifm}
\end{equation}

The work-function difference is not inserted as a single additive voltage
through the whole stack. Instead, the code computes a ferroelectric-specific
built-in voltage partition,
\begin{equation}
\Delta V_{\mathrm{bi,FE}} =
(W_b - W_t)\,
\frac{t_f}{
t_f + \dfrac{k_f}{k_i} t_i + \dfrac{k_f}{k_d} t_d }.
\label{eq:dv-bi-fe-mifm}
\end{equation}
This is the exact expression used in \texttt{src/potential.py}.

Given a polarization $P_{\mathrm{FE}}$ and an applied voltage $V_a$, the active
code computes the screening charge density as
\begin{equation}
\sigma_s(P_{\mathrm{FE}},V_a) =
\frac{
\dfrac{P_d t_d}{k_d}
+
\dfrac{P_{\mathrm{FE}} t_f}{k_f}
+
\epsilon_0 V_a
}{
\dfrac{\lambda_t}{k_t}
+
\dfrac{\lambda_b}{k_b}
+
\dfrac{t_d}{k_d}
+
\dfrac{t_f}{k_f}
+
\dfrac{t_i}{k_i}
}.
\label{eq:sigma-code-mifm}
\end{equation}
The internal field inside the ferroelectric is then
\begin{equation}
E_{\mathrm{FE}}(P_{\mathrm{FE}},V_a) =
\frac{\sigma_s(P_{\mathrm{FE}},V_a)-P_{\mathrm{FE}}}{k_f \epsilon_0}
+
\frac{\Delta V_{\mathrm{bi,FE}}}{t_f}.
\label{eq:efe-code-mifm}
\end{equation}
Equation~\eqref{eq:sigma-code-mifm} says that the electrodes must supply enough
free charge to compensate the external source term, the ferroelectric bound
charge, and any dead-layer contribution while respecting the total
series-capacitance-like denominator of the whole stack. Equation
\eqref{eq:efe-code-mifm} then states that the ferroelectric feels a
screening-controlled electrostatic field together with a built-in contribution
set by the work-function asymmetry and the dielectric partition of the stack.
This is the reason IL thickness, IL dielectric constant, metal work function,
and metal screening length all influence the final hysteresis loop.

\section{Self-consistency condition and bisection solver}

At one applied voltage $V_a$, the solver first freezes the previous domain
population $\{s_n^{\mathrm{base}}\}_{n=1}^{N}$. For a guessed field $E$, it
builds a candidate switched state using
\eqref{eq:switch-rule-mifm} and then computes the corresponding candidate
polarization,
\begin{equation}
P_{\mathrm{FE}}^{\mathrm{cand}}(E) =
\frac{1}{N}\sum_{n=1}^{N} s_n(E) P_{s,n}.
\label{eq:candidate-p-mifm}
\end{equation}
The active bisection solver defines the residual
\begin{equation}
r(E;V_a) =
E - E_{\mathrm{FE}}\!\left(P_{\mathrm{FE}}^{\mathrm{cand}}(E), V_a\right),
\label{eq:residual-mifm}
\end{equation}
and seeks a field $E^\star$ such that
\begin{equation}
r(E^\star;V_a) = 0.
\label{eq:fixed-point-mifm}
\end{equation}
This equation is the appropriate quasistatic closure condition for the present
model because it requires the field used to update the domains to be equal to
the field that the stack electrostatics would actually produce for the
polarization generated by those updated domains.

The active solver first builds an initial bracket by estimating a sufficiently
large field window from the mean magnitude of the local polarization values,
the extreme ferroelectric fields corresponding to $P=\pm P_{\max}$, and the
maximum coercive-field magnitude in the domain population, then adds a margin
of 1\,\mvcm\ in atomic units. If the residual does not change sign across that
interval, the bracket is expanded up to 24 times. This step is purely
numerical; it is not additional physics. The bisection loop itself terminates
when the polarization difference between the high-field and low-field
candidates becomes sufficiently small,
\begin{equation}
\left|
P_{\mathrm{FE}}^{\mathrm{high}} - P_{\mathrm{FE}}^{\mathrm{low}}
\right|
\le \Delta P_{\mathrm{tol}},
\label{eq:ptol-mifm}
\end{equation}
and in the active figure workflow the tolerance is
\begin{equation}
\Delta P_{\mathrm{tol}} = 0.5 \ \uc.
\label{eq:ptol-active-mifm}
\end{equation}
If the bracket cannot be made sign-changing or if the iteration limit is
reached, the solver chooses the candidate with the smallest absolute residual.
That fallback is a numerical protection device rather than a new physical
assumption.

\section{Sweep protocol and active parameter set}

The active figure driver constructs a single-cycle voltage trace
\begin{equation}
-V_{\max} \rightarrow +V_{\max} \rightarrow -V_{\max}
\label{eq:single-cycle-mifm}
\end{equation}
using a fixed voltage step, repeats that cycle twice,
\begin{equation}
\text{full sweep} = \text{cycle} \times 2,
\label{eq:two-cycles-mifm}
\end{equation}
and plots only the final cycle. This reduces the dependence of the displayed
loop on the arbitrary initial domain state. Across the IL-thickness sweep the
same seed and the same structural-distribution parameters are reused, so the
same $c/a$ realization and the same initial random domain-state pattern are
preserved from one IL thickness to the next. This isolates the electrostatic
effect of IL thickness rather than mixing geometry changes with random changes
in the ferroelectric population.

The most important baseline material parameters used by the active figure path
are summarized in Table~\ref{tab:materials-active-mifm}.

\begin{table}[H]
\centering
\caption{Baseline active material parameters used by the current figure workflow.}
\label{tab:materials-active-mifm}
\begin{tabular}{@{}lllll@{}}
\toprule
Material & Work function (eV) & Screening length (bohr) & $k$ & Notes \\
\midrule
Ti & 4.33 & 0.1 & 1 & baseline metal \\
Al & 4.28 & 0.1 & 1 & baseline metal \\
Pd & 5.30 & 0.1 & 1 & baseline metal \\
Test & 4.08 & 1.0 & 1 & lower WF, longer screening \\
Test2 & 4.33 & 1.0 & 1 & same WF as Ti, longer screening \\
Al$_2$O$_3$ & --- & --- & 9.3 & baseline AlO$_x$ proxy \\
HfO$_x$ & --- & --- & 16.64 & high-$k$ interlayer \\
AlScN & --- & --- & 16 & baseline FE dielectric constant \\
\bottomrule
\end{tabular}
\end{table}

The main thickness-specific voltage windows in the active figure set are
45\,nm AlScN at $\pm 20$ V, 20\,nm AlScN at $\pm 12$ V, and 10\,nm AlScN at
$\pm 6.5$ V.

\section{Qualitative consequences of the equations}

Equations~\eqref{eq:sigma-code-mifm} and \eqref{eq:efe-code-mifm} already
predict the dominant trends in the figure set. Increasing the interlayer
thickness raises the electrostatic denominator, reduces the charge-compensation
efficiency of the electrodes, weakens the transmission of external voltage to
the ferroelectric, and strengthens the depolarization penalty; as a result, the
loop generally shrinks within a fixed voltage window. Increasing the
interlayer dielectric constant has the opposite effect in the sense that a
larger $k_i$ reduces the electrostatic penalty associated with the interlayer.
This is why the active model predicts larger loops for HfO$_x$ than for
Al$_2$O$_3$ at the same thickness.

The work-function difference enters through
\eqref{eq:dv-bi-fe-mifm}, so changing the top--bottom order changes the sign of
the built-in term. Reversing a strongly asymmetric stack should therefore
approximately reverse the horizontal preference of the loop. The screening
lengths enter the denominator of \eqref{eq:sigma-code-mifm} directly. Poorer
screening reduces the screening charge available to compensate bound charge,
strengthens the depolarization field, and can shrink both positive and negative
branches even when the work function is unchanged.

\section{Representative active results}

The generated figure set is consistent with these expectations.
Figures~\ref{fig:20nm-al2o3-mifm} and \ref{fig:20nm-hfox-mifm} show a
representative comparison for the 20\,nm AlScN case at $\pm 12$ V.

\begin{figure}[H]
\centering
\includegraphics[width=0.78\linewidth]{20nm_figures/Ti_Al2O3_AlScN_Ti_il_hysteresis_12V.png}
\caption{Representative Al$_2$O$_3$ interlayer sweep: 20\,nm AlScN at $\pm 12$ V.}
\label{fig:20nm-al2o3-mifm}
\end{figure}

\begin{figure}[H]
\centering
\includegraphics[width=0.78\linewidth]{20nm_figures/Ti_HfOx_AlScN_Ti_il_hysteresis_12V.png}
\caption{Representative HfO$_x$ interlayer sweep: 20\,nm AlScN at $\pm 12$ V.}
\label{fig:20nm-hfox-mifm}
\end{figure}

For the 20\,nm case at 3\,nm IL, the Al$_2$O$_3$ stack gives
$P_{\max} = 7.830$ \uc, whereas the HfO$_x$ stack gives
$P_{\max} = 16.533$ \uc. This is the expected electrostatic trend associated
with the larger dielectric constant of HfO$_x$.

\begin{figure}[H]
\centering
\includegraphics[width=0.88\linewidth]{45nm_figures/all_electrodes_0nm_overlay.png}
\caption{Electrode sensitivity at 0\,nm IL.}
\label{fig:electrode-overlay-active-mifm}
\end{figure}

The electrode comparison provides a second useful check. Ti/Test2 differs
strongly from Ti/Ti even though Test2 preserves the Ti work function. This is a
direct signature of screening-length sensitivity in the active electrostatic
model, and it confirms that the metal boundary condition is controlled by
screening as well as work-function asymmetry.

\section{Scope of the active model}

The present model is multidomain, history-dependent, electrostatically
self-consistent, quasistatic, and physically interpretable. It is not a
transport model, a time-domain switching-kinetics model, a continuous Landau or
phase-field model, or a microscopic \textit{ab initio} domain simulation. Those
statements are not incidental; they define the domain of validity of the active
workflow. The repository is solving a compact, interpretable hysteresis problem
rather than a full diode transport problem.

\section{Closing statement}

The active repository solves the following problem: for a fixed multilayer
stack and a fixed Preisach domain population, find the field-dependent domain
configuration whose net polarization is self-consistent with the
one-dimensional stack electrostatics at each point of a quasistatic voltage
sweep. That statement captures the exact physical and mathematical content of
the current multidomain Preisach hysteresis workflow.

\end{document}
